\section{Formal Task and Evidence Audit}
This supplement specifies evidence matching, annotation checks, tool behavior, execution settings and action comparisons.

\subsection{Recorded, obtainable and released evidence}
Let \(q\) be a claim, \(z\) the fixed physical state and \(\mathcal{P}\) the recorded photographic pool.
A configuration specifies the obtainable subset \(\mathcal{A}\subseteq\mathcal{P}\) and initial release \(\mathcal{R}_0\subseteq\mathcal{A}\).
At step \(t\), the public history contains \(q\), released images \(\mathcal{R}_t\), actual tool responses and remaining budget.
A request can change \(\mathcal{R}_t\), but never \(z\). The main paper's \(\mathcal{E}_t\) denotes the displayed pixels of released records and derived views; unseen detail from a private original cannot contribute evidence.

An evidence record contains an opaque image ID, pixels, object slot, source time and camera view.
Private annotations specify condition, identity, visibility, availability and sufficient evidence alternatives.
The policy receives only public observations and permitted provenance.
Source IDs identifying unreleased originals do not authorize citing those originals.

The terminal decision is \(d=(\hat y,\mathcal{C},\mathcal{L})\), where \(\hat y\in\{\Spt,\Rft,\Nme\}\), \(\mathcal{C}\) contains image/region/time citations and \(\mathcal{L}\) contains cross-image relations.
The labels denote \emph{Supported}, \emph{Refuted} and \emph{Need more evidence}.

The minimum acquisition cost is
\begin{equation}
c^\star(q,v)=
\min_{\tau\in\mathcal{T}_v:\,Y_q(\mathcal{E}_{\tau})\ne\Nme}
C(\tau),\qquad \min\varnothing=\infty.
\end{equation}
Here \(\mathcal{T}_v\) contains finite, publicly expressible trajectories obeying access and tool rules before resource caps; \(C\) sums charged costs. Admission additionally requires one sufficient trajectory \(\tau\) satisfying both \(C(\tau)\le B\) and \(\ell(\tau)\le T\), where \(T\) is the turn limit.
All observations needed for a sufficient set must be jointly reachable on one trajectory.
This annotation-side minimum is not a policy observation or a claim to have implemented an exact search solver.
Budget comparisons keep the access target fixed and separately identify cases lacking a joint cost-and-turn-feasible path. The scalar \(c^\star\) alone does not certify turn feasibility.

\subsection{Alternative sufficient evidence sets}
For \(q=\bigwedge_{j=1}^{m}q_j\), each subclaim has truth \(t_j\in\{\Spt,\Rft,?\}\) and alternative sufficient evidence sets \(\mathcal{M}_j^u\), for \(u\in\{\Spt,\Rft\}\).
Each sufficient set specifies required source regions and relations.
Write \(M\preceq d\) when submitted citations and relations satisfy every requirement in \(M\).
Then
\begin{align}
E_{\Spt}(d)
 &=\bigwedge_j\left[t_j=\Spt\ \land\
 \exists M\in\mathcal{M}_j^{\Spt}:M\preceq d\right],\\
E_{\Rft}(d)
 &=\bigvee_j\left[t_j=\Rft\ \land\
 \exists M\in\mathcal{M}_j^{\Rft}:M\preceq d\right].
\end{align}
An uncertain conjunct prevents complete support but does not invalidate a decisive refutation of another conjunct.
Figure~\ref{fig:suppcert} illustrates this rule.

\begin{figure}[t]
\centering
\resizebox{\linewidth}{!}{\begin{tikzpicture}[x=1cm,y=1cm,font=\sffamily\footnotesize,
 cert/.style={draw=teal!65!black,fill=teal!7,rounded corners=2pt,minimum width=1.8cm,minimum height=.72cm,align=center},
 counter/.style={cert,draw=orange!75!black,fill=orange!8},
 arrow/.style={-{Stealth[length=1.5mm]},thick}]
\node[anchor=west,font=\sffamily\bfseries] at (0,3.46) {Example: same chair \(\wedge\) damaged armrest};
\node[cert] (s1) at (.95,2.13) {Identity\\\(M_1^{\Spt}\)};
\node[cert] (s2) at (3.0,2.13) {Condition\\\(M_2^{\Spt}\)};
\node (and) at (4.65,2.13) {AND};
\node[anchor=west] (yes) at (5.65,2.13) {Supported};
\draw[arrow] (s1.north) -- ++(0,.38) -| (and.north);
\draw[arrow] (s2.east) -- (and.west);
\draw[arrow] (and.east) -- (yes.west);
\node[counter] (r1) at (.95,.72) {Different chair\\\(M_1^{\Rft}\)};
\node[counter] (r2) at (3.0,.72) {Armrest intact\\\(M_2^{\Rft}\)};
\node (or) at (4.65,.72) {OR};
\node[anchor=west] (no) at (5.65,.72) {Refuted};
\draw[arrow,draw=orange!75!black] (r1.south) -- ++(0,-.38) -| (or.south);
\draw[arrow,draw=orange!75!black] (r2.east) -- (or.west);
\draw[arrow,draw=orange!75!black] (or.east) -- (no.west);
\node[anchor=west,text=black!65,font=\scriptsize] at (0,-.28) {Each certificate can have alternative sufficient views.};
\end{tikzpicture}
}
\caption{\textbf{Evidence requirements for a conjunction.}
Both identity and condition must be established for support; either reliable counterexample can refute the claim.
This is a logical illustration, without assigned photographic labels.}
\label{fig:suppcert}
\end{figure}

Citations use released image IDs and original-source integer coordinates and times.
Crops preserve a mapping from display coordinates to source coordinates.
The exact union of cited source regions must cover a required region; repeated or overlapping crops cannot inflate coverage.
A full-image citation may satisfy this inclusion check, so it is distinct from localization quality.
Relation endpoints additionally require matching images and times and the annotated IoU threshold (default \(0.5\)).
Same-object, different-object and same-time relations are symmetric; earlier-than is directed.
Any additional incorrect relation invalidates the submitted evidence set.

Let \(Y_q(\mathcal{A})\) be the target justified by the complete obtainable pool: \(\Spt\) when every support requirement is reachable, \(\Rft\) when a decisive counterexample is reachable, and \(\Nme\) otherwise.
Evidence validity is checked using only released evidence.
This separates incomplete acquisition from a genuinely unresolved case.

\subsection{Semantic admission checklist}
Two annotators independently assess claim meaning, object identity, condition, visibility, relevant times and sufficient evidence alternatives.
They inspect both initial and obtainable material for every access configuration.
Adjudication must establish a justified target or leave the case pending. Pending annotation status is not the task label \(\Nme\). Review uses the exposed displays, registered resolution and permitted crop paths, not private high-resolution files alone.
An intended Obtainable case must have a sufficient path within the declared action and turn budget, using requests expressible through public object, time and view fields. The evaluator's knowledge of a private filename cannot supply that path.
Preview resolution and a failed request cannot by themselves determine a semantic verdict.

The 64 source controls (32 per source label) and 226 identity candidates remain unreviewed; neither source labels nor missing truth determine a semantic target.

Related views, rewrites and release versions remain grouped.
Families sharing an image or a CO3D sequence form connected source groups for splitting and uncertainty estimation.
Each access version starts a fresh conversation, without history or memory from sibling versions. The pool is evaluation-only; future training splits must be fixed by source group before fitting a policy.

The separate interface study targets 20 participants, each completing three search trials and two full-pool judgments in 5--10 minutes, without hard time or action limits. It checks clarity; optional region marks and full-pool judgments do not replace evidence annotation.

\section{Environment and Tool Contract}
\subsection{Observations, costs and provenance}
The tool pool and default costs are shown in Table~\ref{tab:supptools}.
Failed requests incur their ordinary request cost.
Metadata timestamps are submitted provenance claims unless independently corroborated.

\begin{table}[t]
\centering\small
\caption{\textbf{Default tool costs and observations.}}
\label{tab:supptools}
\begin{tabularx}{\linewidth}{@{}l r Y@{}}
\toprule
Action & Cost & Observation\\
\midrule
Inspect & 1 & Released image or derived view\\
Crop / zoom & 1 & Resampled source pixels with mapping\\
OCR & 2 & Text, polygons and confidence\\
Compare & 2 & Two released pixel views\\
Request photo & 3 & Matching record or failure\\
Image quality & 1 & Pixel-derived quality statistics\\
Metadata & 1 & Permitted provenance and file hash\\
Text-to-image & 3 & Grounding DINO regions\\
Image-to-image & 3 & DINOv2 matching proposals\\
Image-to-text & 2 & CLIP candidate ranking\\
Finish & 0 & Verdict, citations and relations\\
\bottomrule
\end{tabularx}
\end{table}

Crop inputs use display coordinates; finish citations use original-source coordinates.
Derived views retain source rectangle, original dimensions, display dimensions, time and operation history.
Enlarging a downsampled view cannot recover discarded detail.
A derived view is not independent evidence.

Grounding DINO proposes text-conditioned regions~\cite{groundingdino}; DINOv2 features propose image correspondences~\cite{dinov2}; CLIP ranks user-supplied descriptions~\cite{clip}.
Similarity and confidence scores are not identity truth or calibrated verdict probabilities.
An empty grounding result does not prove absence.
OCR outputs remain fallible measurements.

Requests match the public object slot, time and view against admissible records, rather than retrieving the most informative record automatically.
Missing and Unavailable use equivalent failure semantics when neither can return a photograph.
Corrective errors disclose only released IDs, permitted bounds and public action constraints.
They reveal no private sufficiency score or unreleased asset.

\subsection{Isolation and simulation}
The policy worker receives public messages and tool schemas.
A private controller validates actions, executes them and audits terminal evidence.
The policy-facing interface exposes no annotation, evaluation or fork endpoint.
Policies allowed arbitrary code execution require filesystem or container isolation; process separation alone is not a filesystem barrier.

The current public adapter releases recorded photographs from a closed pool.
It cannot capture an unrecorded view.
State-driven rendering is a separate simulation capability; its integration fixtures are not real-photo evaluation cases.
Generated enhancement and imagined photographs cannot become supporting evidence.
A future generative world model could propose a viewing direction while keeping its predictions separate from source observations.

\section{Execution and Scoring Details}
\subsection{Initial evaluation configuration}
The initial local model matrix comprises Qwen2.5-VL-7B, Qwen3-VL-8B, Qwen3.5-9B and InternVL3.5-8B~\cite{qwen25,qwen3,qwen35,internvl35}.
The configuration uses BF16 weights, greedy generation, a tool budget of 12, at most eight assistant turns and an 8,192-token context cap.
Direct decisions permit at most 96 output tokens; interactive turns permit 384.
Input images are capped at \(512\times512\) pixels, without additional InternVL tiling.
Checkpoint-native preprocessing and any architecture-specific behavior remain recorded in manifests.

Direct and interactive output caps differ, so their score gap cannot isolate acquisition. Capability attribution requires replaying final decisions with matched output allowance and formatting, auditing presented images and resolution, and reporting truncation separately. This control has not been run.

Qwen3.5 uses a separate Transformers 5.3.0 environment with thinking disabled.
The other three models use Transformers 4.57.6.
Model/data revisions, source digests, providers and dependency versions are recorded per run.
Runs with different interface diagnostics or source settings remain separate.

Native JSON/XML and portable structured calls map to common tool schemas.
Adapters do not invent arguments, repair the selected function or accept multiple calls as a single valid call.
Checkpoint stopping tokens and the assistant-turn boundary are preserved so control returns to the real environment.
The next observation must come from actual execution, never a model-generated tool response.

\subsection{Metrics and failure accounting}
For reviewed case \(i\), let \(Y_i=Y_{q_i}(\mathcal{A}_i)\) and \(K_i=1\) when a definite verdict has a valid evidence set citing only released evidence.
For reviewed families, strict family success is
\begin{equation}
\mathrm{FG}=\frac{1}{|\mathcal{F}|}
\sum_{f\in\mathcal{F}}\prod_{v\in\mathcal{V}_f}G(d_{fv}),
\end{equation}
where \(\mathcal{V}_f\) contains four admitted access versions.
It complements average accuracy; access versions are not independent samples.
The main paper's acquisition--conversion identity uses only reviewed answerable episodes with a jointly cost-and-turn-feasible sufficient path. Each attempted episode enters its denominator, including failed execution. Evidence exposure is audited before termination regardless of the submitted verdict; exposed pixels must satisfy the same reviewed sufficiency requirements used to check citations. Conversion includes perception, reasoning, citation and output-format failures, so it is not a pure perception score.

Grounded accuracy and unsupported decision rate are
\begin{align}
\mathrm{GA}
 &=\frac{1}{n}\sum_i\ind[\hat y_i=Y_i]\,
   \ind[\hat y_i=\Nme\ \lor\ K_i=1],\\
\mathrm{UDR}
 &=\frac{1}{n}\sum_i
   \ind[\hat y_i\in\{\Spt,\Rft\}\land K_i=0].
\end{align}
Unsubmitted verdicts contribute zero to grounded accuracy.
They are never converted into \(\Nme\).

Coverage is the fraction of cases receiving a definite verdict.
Its conditional risk is reported jointly with coverage; risk is undefined at zero coverage.
A continuous coverage--risk curve requires a declared confidence score and threshold rule~\cite{selective}.
Categorical finish labels alone do not supply calibrated confidence.

Reporting separates valid decisions from parse failure, execution failure and context, turn or budget exhaustion.
Request count, tool cost, tokens, model latency and episode latency are logged separately.
Correspondence accuracy and damage-region IoU apply only to suitably annotated subsets.
Until semantic admission is complete, candidate-level reports are limited to interaction behavior and explicitly labeled construction-label agreement.

\section{Controlled Action Comparisons}
A trusted evaluator forks the same public conversation, released pool, derived views and remaining budget while preserving the physical world.
The candidate action is executed and charged before continuation.
All branches use the same continuation policy and total remaining-turn allowance.
Stochastic continuations additionally require matched recorded seeds.

The audit records evidence released by the action and continuation, final citations and cost. Acquisition succeeds when sufficient evidence becomes available; grounding requires valid terminal citations. Pixel dependence requires the content controls below.
The stop branch asks the same policy for an immediate terminal decision from the shared history; it is never assigned a gold verdict. Final-decision formatting and output allowance are matched across branches.

Attributing benefit to returned content requires additional controls.
Matched random patches are relevant for crop actions~\cite{countercredit}.
Comparisons involving recorded views must also account for source, view type, resolution and cost.
Without a defensible content control, the analysis reports an action intervention rather than a causal attribution to returned pixels.

An obtainable and unavailable configuration can share an identical public prefix.
Before access feedback is observed, hidden knowledge of the configuration cannot justify different preferred actions.
Action-value estimates must aggregate over matched histories or condition only on visible feedback.

Future training exports retain reviewed assistant actions and actual observations, excluding private labels from prompts.
Sibling comparisons must share the public prefix, schemas and world state.
Unreviewed or proxy traces are excluded; no training gain or new optimizer is claimed.
